% Generated draft: requires mathematical and source review.
\documentclass[11pt]{article}
\usepackage[margin=25mm]{geometry}
\usepackage{fontspec}
\setmainfont{Latin Modern Roman}
\usepackage{amsmath,amssymb,mathtools,unicode-math}
\setmathfont{Latin Modern Math}
\usepackage{enumitem,needspace,xurl,hyperref}
\setlength{\parindent}{0pt}
\setlength{\parskip}{5pt}
\setlength{\emergencystretch}{3em}
\newcommand{\N}{\mathbb{N}}
\newcommand{\Z}{\mathbb{Z}}
\newcommand{\Q}{\mathbb{Q}}
\newcommand{\R}{\mathbb{R}}
\newcommand{\C}{\mathbb{C}}
\newcommand{\F}{\mathbb{F}}
\newcommand{\A}{\mathbb{A}}
\newcommand{\PP}{\mathbb{P}}
\newcommand{\E}{\mathbb{E}}
\newcommand{\eps}{\varepsilon}
\newcommand{\dd}{\,\mathrm{d}}
\begin{document}
\section*{UnsolvedMath ID: 30003948}
\textbf{Source problem code:} catalogue entry
\subsection{Definitions and mathematical statement}
Let \(S\) be the compact hyperbolic surface in the supplied record. Write \(\phi_n\) for its Laplace eigenfunctions and \(\lambda_n\) for the corresponding eigenvalues. With the nonnegative sign convention for the Laplacian, the eigenvalue equation is \(\Delta_S\phi_n=\lambda_n\phi_n\); the source does not explicitly specify this sign convention, an eigenfunction normalization, or boundary conditions if \(S\) has boundary.

For a simple closed geodesic \(C\subset S\), let \(\ell(C)\) denote its hyperbolic length and let \(\mathrm{d}s\) denote hyperbolic arclength. Put
\[
\mathcal G_S(T)=\{C\subset S:C\text{ is a simple closed geodesic and }\ell(C)\leq T\},\qquad N_S(T)=\#\mathcal G_S(T).
\]
Here the set notation treats geodesics as geometric curves counted once; the supplied record does not explicitly state its counting convention. Whenever \(0<N_S(T)<\infty\), define the signed average of geodesic periods by
\[
A_n(T)=\frac{1}{N_S(T)}\sum_{C\in\mathcal G_S(T)}\int_C\phi_n\,\mathrm{d}s.
\]
The integrals are not divided by \(\ell(C)\), and no absolute values are placed inside the sum.

The supplied question asks whether there exist constants \(\alpha,\beta>0\) such that
\[
A_n(T)=O\!\left(\lambda_n^{-\alpha}T^{-\beta}\right).
\]
For positive eigenvalues and positive length cutoffs with nonzero denominator, an estimate of this form would mean
\[
\left|\sum_{C\in\mathcal G_S(T)}\int_C\phi_n\,\mathrm{d}s\right|
\leq K\,N_S(T)\lambda_n^{-\alpha}T^{-\beta}
\]
with \(K>0\) independent of the variables over which the estimate is asserted to be uniform. The source does not specify that uniformity, the joint asymptotic regime in \(n\) and \(T\), or any relation between these variables. It also does not specify whether the exponents and implicit constant may depend on \(S\), the eigenbasis, or the individual eigenfunction. The expression with \(\lambda_n^{-\alpha}\) requires a restriction to positive eigenvalues or a separately stated treatment of zero eigenvalues, which is absent. These qualifications remain unresolved; the displayed target is therefore a statement-only transcription requiring clarification, not a fully specified conjecture.
\subsection{Short English statement}
The supplied problem asks whether averaging the signed integrals of a Laplace eigenfunction over simple closed geodesics of length at most \(T\) gives simultaneous power decay in its eigenvalue and in \(T\). Eigenfunction normalization, the treatment of zero eigenvalues, surface boundary assumptions, and the intended uniformity and limiting regime require clarification.
\subsection{Sources}
\begin{itemize}
\item Supplied local UnsolvedMath registry record, permanent ID \(30003948\), title: \emph{Eigenfunction Averages Along Simple Closed Geodesics}; source set: Oberwolfach Reports Open Problems; exact source problem code: \texttt{catalogue entry}. The record gives proposed year \(2018\), leaves the proposer unspecified, and supplies the labels \texttt{open} and \texttt{corrected\_verified}. These are registry labels, not independent mathematical verification.
\item Supplied report citation: \emph{New Trends in Teichmüller Theory and Mapping Class Groups} (2019), Oberwolfach Reports, pp.\ 2475--2534, DOI \texttt{10.4171/owr/2018/40}. \url{https://doi.org/10.4171/owr/2018/40}. The supplied literature-source label is ``Oberwolfach Report 40, Burger problem.'' That label is retained as supplied attribution; it does not resolve the empty proposer field.
\item Supplied literature reference: Reznikov, geodesic restrictions of eigenfunctions. \url{https://arxiv.org/abs/math/0403437}.
\item Supplied literature reference: Sogge--Xi--Zhang, geodesic period integrals. \url{https://arxiv.org/abs/1604.03189}.
\item Literature assessment dated \texttt{2026-08-21}: the supplied curation labels the problem open and states that known restriction estimates for eigenfunctions on individual geodesics do not prove the requested averaged estimate with simultaneous eigenvalue and length decay. This is supplied dated curation derived from the Oberwolfach Reports material, not an independently verified literature assessment. The supplied background explicitly requests independent verification of time-sensitive status and cited claims.
\end{itemize}
\end{document}
